\documentclass[10pt,a4paper]{amsart}
\usepackage{amsmath,amssymb,amsthm,graphicx}
% Retained preamble for the Stacks Project chapter snapshot in this directory.
% The chapter's own preamble.tex is a copy of the shared Stacks Project preamble; the block
% below is sliced MECHANICALLY, line for line, from it. Nothing was added, removed or
% redefined; the wrapper class, the xr-hyper/hyperref wiring and the externaldocument list
% are omitted because the assembler supplies the document class.
%
% lines taken from preamble.tex: 3, 5, 9, 11, 12, 156, 157, 158, ... (51 lines total)
%
% package coverage restored from the chapter's own preamble (2026-10-13): xy, xy-2cell

\usepackage[all]{xy}
\xyoption{2cell}
\UseAllTwocells

\newenvironment{reference}{\comment}{\endcomment}
\newenvironment{slogan}{\comment}{\endcomment}
\newenvironment{history}{\comment}{\endcomment}
\theoremstyle{plain}
\newtheorem{theorem}[subsection]{Theorem}
\newtheorem{proposition}[subsection]{Proposition}
\newtheorem{lemma}[subsection]{Lemma}
\theoremstyle{definition}
\newtheorem{definition}[subsection]{Definition}
\newtheorem{example}[subsection]{Example}
\newtheorem{exercise}[subsection]{Exercise}
\newtheorem{situation}[subsection]{Situation}
\theoremstyle{remark}
\newtheorem{remark}[subsection]{Remark}
\newtheorem{remarks}[subsection]{Remarks}
\numberwithin{equation}{subsection}
\def\lim{\mathop{\mathrm{lim}}\nolimits}
\def\colim{\mathop{\mathrm{colim}}\nolimits}
\def\Spec{\mathop{\mathrm{Spec}}}
\def\Hom{\mathop{\mathrm{Hom}}\nolimits}
\def\Ext{\mathop{\mathrm{Ext}}\nolimits}
\def\SheafHom{\mathop{\mathcal{H}\!\mathit{om}}\nolimits}
\def\SheafExt{\mathop{\mathcal{E}\!\mathit{xt}}\nolimits}
\def\Sch{\mathit{Sch}}
\def\Mor{\mathop{\mathrm{Mor}}\nolimits}
\def\Ob{\mathop{\mathrm{Ob}}\nolimits}
\def\Sh{\mathop{\mathit{Sh}}\nolimits}
\def\NL{\mathop{N\!L}\nolimits}
\def\CH{\mathop{\mathrm{CH}}\nolimits}
\def\proetale{{pro\text{-}\acute{e}tale}}
\def\etale{{\acute{e}tale}}
\def\QCoh{\mathit{QCoh}}
\def\Ker{\mathop{\mathrm{Ker}}}
\def\Im{\mathop{\mathrm{Im}}}
\def\Coker{\mathop{\mathrm{Coker}}}
\def\Coim{\mathop{\mathrm{Coim}}}
\DeclareMathSymbol{\boxtimes}{\mathbin}{AMSa}{"02}
\def\QCohstack{\mathcal{QC}\!\mathit{oh}}
\def\Cohstack{\mathcal{C}\!\mathit{oh}}
\def\Spacesstack{\mathcal{S}\!\mathit{paces}}
\def\Quotfunctor{\mathrm{Quot}}
\def\Hilbfunctor{\mathrm{Hilb}}
\def\Curvesstack{\mathcal{C}\!\mathit{urves}}
\def\Polarizedstack{\mathcal{P}\!\mathit{olarized}}
\def\Complexesstack{\mathcal{C}\!\mathit{omplexes}}
\def\Pic{\mathop{\mathrm{Pic}}\nolimits}
\def\Picardstack{\mathcal{P}\!\mathit{ic}}
\def\Picardfunctor{\mathrm{Pic}}
\def\Deformationcategory{\mathcal{D}\!\mathit{ef}}
\usepackage{cleveref}
\crefname{originallemma}{original lemma}{original lemmas}
\Crefname{originallemma}{Original lemma}{Original lemmas}
\makeatletter
\newlabel{algebra-lemma-CM-over-regular-flat}{{00R4}{source}{}{sourcepointer-1}{}}
\newlabel{algebra-lemma-CM-over-regular-flat@cref}{{[originallemma][0][]00R4}{[1][0][]source}{}{}{}}
\makeatother
\crefname{originallemma}{original lemma}{original lemmas}
\Crefname{originallemma}{Original lemma}{Original lemmas}
\makeatletter
\newlabel{divisors-lemma-characterize-OD}{{01X0}{source}{}{sourcepointer-2}{}}
\newlabel{divisors-lemma-characterize-OD@cref}{{[originallemma][0][]01X0}{[1][0][]source}{}{}{}}
\makeatother
\crefname{originallemma}{original lemma}{original lemmas}
\Crefname{originallemma}{Original lemma}{Original lemmas}
\makeatletter
\newlabel{lemma-different-syntomic-quasi-finite}{{0BWG}{source}{}{sourcepointer-3}{}}
\newlabel{lemma-different-syntomic-quasi-finite@cref}{{[originallemma][0][]0BWG}{[1][0][]source}{}{}{}}
\makeatother
\crefname{originallemma}{original lemma}{original lemmas}
\Crefname{originallemma}{Original lemma}{Original lemmas}
\makeatletter
\newlabel{more-morphisms-lemma-lci-permanence}{{069M}{source}{}{sourcepointer-4}{}}
\newlabel{more-morphisms-lemma-lci-permanence@cref}{{[originallemma][0][]069M}{[1][0][]source}{}{}{}}
\makeatother
\crefname{originallemma}{original lemma}{original lemmas}
\Crefname{originallemma}{Original lemma}{Original lemmas}
\makeatletter
\newlabel{more-morphisms-lemma-morphism-between-flat-Noetherian}{{039D}{source}{}{sourcepointer-5}{}}
\newlabel{more-morphisms-lemma-morphism-between-flat-Noetherian@cref}{{[originallemma][0][]039D}{[1][0][]source}{}{}{}}
\makeatother
\crefname{originallemma}{original lemma}{original lemmas}
\Crefname{originallemma}{Original lemma}{Original lemmas}
\makeatletter
\newlabel{morphisms-lemma-smooth-omega-finite-locally-free}{{02G1}{source}{}{sourcepointer-6}{}}
\newlabel{morphisms-lemma-smooth-omega-finite-locally-free@cref}{{[originallemma][0][]02G1}{[1][0][]source}{}{}{}}
\makeatother
\crefname{originallemma}{original lemma}{original lemmas}
\Crefname{originallemma}{Original lemma}{Original lemmas}
\makeatletter
\newlabel{morphisms-lemma-triangle-differentials}{{01UX}{source}{}{sourcepointer-7}{}}
\newlabel{morphisms-lemma-triangle-differentials@cref}{{[originallemma][0][]01UX}{[1][0][]source}{}{}{}}
\makeatother
\title{Stacks, Discriminants: lemma-discriminant-quasi-finite-morphism-smooth}
\author{The Stacks Project Authors}
\date{stacks-project:discriminant.tex:local-snapshot-20261012; GNU FDL 1.2 or later}
\begin{document}
\maketitle
\noindent\textit{Editorial note.} Statement span is lines 2765-2785 of sources/stacks-discriminant/pinned-original-source.tex; proof span is lines 2787-2814, the complete author proof (26 body lines).\par\medskip
\noindent\textit{Editorial note.} Navigation stubs: 7 cited-but-not-included labels are supplied as non-load-bearing pointers with their pinned tags where the snapshot tag index has them.\par\medskip
\noindent\textit{Editorial note.} extra\_packages is empty on purpose: the retained preamble already loads xy with options and a second declaration clashes in the assembler wrapper.\par\medskip
\noindent\textit{Editorial note.} Text is retained verbatim, including the author's own typos and placeholders.\par\medskip
\noindent\textit{Editorial note.} Queue index 17 of queue\_discriminant.json (pb=1352, stmt 2765-2785, proof 2787-2814).\par\medskip
\noindent\textit{Allowed standard prerequisites (editorial).}\begin{itemize}
\item Explicit prerequisite boundary. The cited inputs are prerequisites; the argument in the retained span is not replaced by any cited result.
\end{itemize}\medskip
\ifdefined\hypertarget\hypertarget{sourcepointer-1}{}\fi
\noindent\textit{Source reference.} Original lemma 00R4: sources/stacks-discriminant/pinned-original-source.tex, source lines 2814--2814. Cited by the retained statement or proof but not part of the collected spans; supplied only as a navigation pointer. The argument does not depend on its content.\par\medskip
\ifdefined\hypertarget\hypertarget{sourcepointer-2}{}\fi
\noindent\textit{Source reference.} Original lemma 01X0: sources/stacks-discriminant/pinned-original-source.tex, source lines 2814--2814. Cited by the retained statement or proof but not part of the collected spans; supplied only as a navigation pointer. The argument does not depend on its content.\par\medskip
\ifdefined\hypertarget\hypertarget{sourcepointer-3}{}\fi
\noindent\textit{Source reference.} Original lemma 0BWG: sources/stacks-discriminant/pinned-original-source.tex, source lines 2708--2713. Cited by the retained statement or proof but not part of the collected spans; supplied only as a navigation pointer. The argument does not depend on its content.\par\medskip
\ifdefined\hypertarget\hypertarget{sourcepointer-4}{}\fi
\noindent\textit{Source reference.} Original lemma 069M: sources/stacks-discriminant/pinned-original-source.tex, source lines 2814--2814. Cited by the retained statement or proof but not part of the collected spans; supplied only as a navigation pointer. The argument does not depend on its content.\par\medskip
\ifdefined\hypertarget\hypertarget{sourcepointer-5}{}\fi
\noindent\textit{Source reference.} Original lemma 039D: sources/stacks-discriminant/pinned-original-source.tex, source lines 2814--2814. Cited by the retained statement or proof but not part of the collected spans; supplied only as a navigation pointer. The argument does not depend on its content.\par\medskip
\ifdefined\hypertarget\hypertarget{sourcepointer-6}{}\fi
\noindent\textit{Source reference.} Original lemma 02G1: sources/stacks-discriminant/pinned-original-source.tex, source lines 2814--2814. Cited by the retained statement or proof but not part of the collected spans; supplied only as a navigation pointer. The argument does not depend on its content.\par\medskip
\ifdefined\hypertarget\hypertarget{sourcepointer-7}{}\fi
\noindent\textit{Source reference.} Original lemma 01UX: sources/stacks-discriminant/pinned-original-source.tex, source lines 2814--2814. Cited by the retained statement or proof but not part of the collected spans; supplied only as a navigation pointer. The argument does not depend on its content.\par\medskip
\begin{lemma}
\label{lemma-discriminant-quasi-finite-morphism-smooth}
Let $S$ be a Noetherian scheme. Let $X$, $Y$ be smooth schemes
of relative dimension $n$ over $S$. Let $f : Y \to X$ be a
locally quasi-finite morphism over $S$.
Then $f$ is flat and the closed subscheme $R \subset Y$
cut out by the different of $f$ is the locally principal
closed subscheme cut out by
$$
\wedge^n(\text{d}f) \in
\Gamma(Y,
(f^*\Omega^n_{X/S})^{\otimes -1} \otimes_{\mathcal{O}_Y} \Omega^n_{Y/S})
$$
If $f$ is \'etale at the associated points of $Y$, then $R$ is an
effective Cartier divisor and
$$
f^*\Omega^n_{X/S} \otimes_{\mathcal{O}_Y} \mathcal{O}(R) =
\Omega^n_{Y/S}
$$
as invertible sheaves on $Y$.
\end{lemma}

\begin{proof}
To prove that $f$ is flat, it suffices to prove $Y_s \to X_s$
is flat for all $s \in S$ (More on Morphisms, Lemma
\ref{more-morphisms-lemma-morphism-between-flat-Noetherian}).
Flatness of $Y_s \to X_s$ follows from
Algebra, Lemma \ref{algebra-lemma-CM-over-regular-flat}.
By More on Morphisms, Lemma
\ref{more-morphisms-lemma-lci-permanence}
the morphism $f$ is a local complete intersection morphism.
Thus the statement on the different follows from the
corresponding statement on the K\"ahler different by
Lemma \ref{lemma-different-syntomic-quasi-finite}.
Finally, since we have the exact sequence
$$
f^*\Omega_{X/S} \xrightarrow{\text{d}f} \Omega_{Y/S} \to \Omega_{Y/X} \to 0
$$
by Morphisms, Lemma \ref{morphisms-lemma-triangle-differentials}
and since $\Omega_{X/S}$ and $\Omega_{Y/S}$ are finite locally free
of rank $n$ (Morphisms, Lemma
\ref{morphisms-lemma-smooth-omega-finite-locally-free}),
the statement for the K\"ahler different is clear from the definition
of the zeroth fitting ideal. If $f$ is \'etale at the associated
points of $Y$, then $\wedge^n\text{d}f$ does not vanish in
the associated points of $Y$, which implies that the local equation
of $R$ is a nonzerodivisor. Hence $R$ is an effective Cartier divisor.
The canonical isomorphism sends $1$ to $\wedge^n\text{d}f$, see
Divisors, Lemma \ref{divisors-lemma-characterize-OD}.
\end{proof}

\end{document}
